% Chapter 17: E-Value: Evidence for Causation in Observational Studies with Unmeasured Confounding
% A First Course in Causal Inference - Peng Ding
% Rewrite V2 - After Team Lead review

\chapter{E-Value：未测混杂观测性研究中的因果证据}\label{ch:17}

\section{引言：当可忽略性假设失效时}\label{sec:17-intro}

\textit{``It is possible that the entire relationship between cigarette smoking and lung cancer is due to some common cause.''}
--- \citep{Fisher1957}

在 Part III 中，我们详细讨论了观测性研究的四大支柱：outcome regression（第 \ref{ch:10} 章）、逆倾向得分加权（第 \ref{ch:11} 章）、双重稳健（第 \ref{ch:12} 章）以及匹配（第 \ref{ch:15} 章）。这四种方法的核心假设是\textbf{可忽略性（ignorability）}：
\begin{equation}
Z \perp \{Y(1), Y(0)\} \mid X .
\label{eq:17-ignorability}
\end{equation}

即在协变量 $X$ 条件下，治疗分配 $Z$ 与潜在结果相互独立。这个假设保证了我们可以从观测数据中识别因果效应。然而，\textbf{可忽略性假设无法从数据本身验证}——它是一种根本性的假设，依赖于研究者的领域知识和对因果结构的判断。

这正是观测性研究长期被批评的原因：如果无法控制所有混淆因素，观察到的治疗与结果之间的关联可能是虚假的。统计学能告诉我们观测数据中的关联有多强，但无法保证这种关联背后存在真实的因果关系。

然而，\textbf{并非所有观测性研究都生来平等}。一些研究即便存在未测混杂，其提供的因果证据也比其他研究更强。关键在于：要完全解释掉观察到的关联，未测混杂需要有多强？如果这种强度在领域知识上是不合理的，我们就有理由相信因果关系确实存在。

本章引入的 \textbf{E-value}（VanderWeele and Ding, 2017）正是对这一思想的数学量化。它基于 Ding and VanderWeele (2016) 的理论，为评估观测性研究中的因果证据提供了一个统一且可解释的度量。第 \ref{ch:18} 章讨论平均因果效应的灵敏度分析；第 \ref{ch:19} 章讨论 Rosenbaum 框架下的匹配研究灵敏度分析。

\section{Cornfield 型灵敏度分析}\label{sec:17-cornfield}

\subsection{从可忽略性到潜在可忽略性}\label{sec:17-latent-ignorability}

放弃假设 \cref{eq:17-ignorability} 并不意味着放弃一切分析工具。我们仍然可以假设\textbf{潜在可忽略性（latent ignorability）}：
\begin{equation}
Z \perp \{Y(1), Y(0)\} \mid (X, U) ,
\label{eq:17-latent-ignorability}
\end{equation}

其中 $U$ 是未测混杂变量，条件于 $(X,U)$ 时治疗与潜在结果独立。注意 \cref{eq:17-latent-ignorability} 比 \cref{eq:17-ignorability} 更弱——它允许 $Z$ 与 $Y$ 之间存在通过 $U$ 传递的关联，只要这种关联可以被 $U$ 完全解释。

本章的方法在二值结果 $Y$ 上表现最佳（Ding and VanderWeele, 2016），但可推广到其他非负结果。为简化叙述，除非特别说明，我们均假设 $Y$ 为二值。

\subsection{真实风险比与观测风险比}\label{sec:17-true-vs-observed-rr}

在 $X=x$ 条件下，真实的条件风险比为：
\begin{equation}
\RR_{ZY\mid x}^{\mathrm{true}} = \frac{\pr\{Y(1)=1 \mid X=x\}}{\pr\{Y(0)=1 \mid X=x\}} .
\label{eq:17-true-rr}
\end{equation}

但我们从数据中能直接计算的是观测风险比：
\begin{equation}
\RR_{ZY\mid x}^{\mathrm{obs}} = \frac{\pr(Y=1 \mid Z=1, X=x)}{\pr(Y=1 \mid Z=0, X=x)} .
\label{eq:17-obs-rr}
\end{equation}

在存在未测混杂 $U$ 时，二者一般不相等。直观上，这是因为 $U$ 在治疗组和对照组中的分布不同（$U$ 同时影响治疗分配和结果），而真实的因果效应需要控制 $U$ 后才能得到。

考虑如下因果图（条件于 $X$）：

\[
Z \leftarrow U \rightarrow Y ,
\]

其中 $U$ 是 $Z$ 和 $Y$ 的共同原因（common cause）。在此结构下，$Z \perp Y \mid (X,U)$ 但 $Z \not\perp Y \mid X$。为简化，我们假设 $U$ 为二值，尽管 Ding and VanderWeele (2016) 的理论对 $U$ 的一般形式也成立。

定义两个灵敏度参数：
\begin{equation}
\RR_{ZU\mid x} = \frac{\pr(U=1 \mid Z=1, X=x)}{\pr(U=1 \mid Z=0, X=x)} ,
\label{eq:17-sensitivity-zu}
\end{equation}

衡量\textbf{治疗-混杂关联}；以及
\begin{equation}
\RR_{UY\mid x} = \frac{\pr(Y=1 \mid U=1, X=x)}{\pr(Y=1 \mid U=0, X=x)} ,
\label{eq:17-sensitivity-uy}
\end{equation}

衡量\textbf{混杂-结果关联}。这两个量都无法从数据中识别——它们是未知的灵敏度参数，反映了未测混杂的强度。

\subsection{Cornfield 不等式}\label{sec:17-cornfield-inequality}

下面的定理是整个 E-value 理论的基石。

\begin{Theorem}[Cornfield-type 上界]\label{th:17-cornfield}
在假设 \cref{eq:17-latent-ignorability} 下，假设
\begin{equation}
\RR_{ZY\mid x}^{\mathrm{obs}} > 1, \quad \RR_{ZU\mid x} > 1, \quad \RR_{UY\mid x} > 1 .
\label{eq:17-assumptions}
\end{equation}
则有
\begin{equation}
\RR_{ZY\mid x}^{\mathrm{obs}} \leq \frac{\RR_{ZU\mid x} \cdot \RR_{UY\mid x}}{\RR_{ZU\mid x} + \RR_{UY\mid x} - 1} .
\label{eq:17-cornfield-bound}
\end{equation}
\end{Theorem}

\textbf{定理 \ref{th:17-cornfield} 的含义}：在潜在可忽略性下，治疗与结果之间的观测关联完全来自两条路径——$Z \leftarrow U \rightarrow Y$——的组合效应。上界 $\RR_{ZU\mid x} \cdot \RR_{UY\mid x} / (\RR_{ZU\mid x} + \RR_{UY\mid x} - 1)$ 体现了这两条路径的"乘法阻塞"：要产生一个较大的 $\RR_{ZY\mid x}^{\mathrm{obs}}$，两条路径都必须足够强。

特别地，从 \cref{eq:17-cornfield-bound} 可推出 \textbf{Cornfield 不等式}（Gastwirth et al., 1998）：
\begin{equation}
\min(\RR_{ZU\mid x}, \RR_{UY\mid x}) \geq \RR_{ZY\mid x}^{\mathrm{obs}} .
\label{eq:17-cornfield-inequality}
\end{equation}

即：要完全解释掉观测到的风险比，未测混杂与治疗的关联以及与结果的关联，都必须\textbf{至少与观测到的风险比一样强}。这个简单而强大的结论，首次由 Cornfield et al. (1959) 在讨论吸烟与肺癌时得到。

\textbf{证明思路}：定义 $f_{1,x} = \pr(U=1 \mid Z=1, X=x)$ 和 $f_{0,x} = \pr(U=1 \mid Z=0, X=x)$。通过全概率公式分解 $\RR_{ZY\mid x}^{\mathrm{obs}}$ 并利用 $Z \perp Y \mid (X,U)$，可得恒等式 (\ref{eq:17-rr-identity})。\footnote{完整证明推导见附录 \cref{sec:appendix-17-proof-theorem}}

\section{E-value}\label{sec:17-evalue}

\subsection{从 Cornfield 不等式到 E-value}\label{sec:17-from-cornfield-to-evalue}

从定理 \ref{th:17-cornfield} 和引理 \ref{lm:17-beta-properties}(4)，定义 $w = \max(\RR_{ZU\mid x}, \RR_{UY\mid x})$，则
\begin{equation}
\RR_{ZY\mid x}^{\mathrm{obs}} \leq \frac{w^2}{2w - 1} .
\label{eq:17-w-bound}
\end{equation}

解这个二次不等式，较小根始终 $\leq 1$，故
\begin{equation}
w \geq \RR_{ZY\mid x}^{\mathrm{obs}} + \sqrt{\RR_{ZY\mid x}^{\mathrm{obs}}(\RR_{ZY\mid x}^{\mathrm{obs}} - 1)} .
\label{eq:17-w-lower-bound}
\end{equation}

这给出了 E-value 的定义。

\begin{Definition}[E-value]\label{def:17-evalue}
给定观测条件风险比 $\RR_{ZY\mid x}^{\mathrm{obs}}$，定义 \textbf{E-value} 为
\begin{equation}
\mathrm{E\text{-}value} = \RR_{ZY\mid x}^{\mathrm{obs}} + \sqrt{\RR_{ZY\mid x}^{\mathrm{obs}}(\RR_{ZY\mid x}^{\mathrm{obs}} - 1)} .
\label{eq:17-evalue}
\end{equation}
\end{Definition}

\textbf{E-value 的解释}：E-value 表示要完全解释掉观测到的风险比，$\max(\RR_{ZU\mid x}, \RR_{UY\mid x})$ 必须达到的最小值。换言之，如果领域知识认为任何合理的未测混杂，其与治疗和结果的关联程度都不可能超过某个阈值 $R$，那么当 $R < \mathrm{E\text{-}value}$ 时，我们有理由相信因果关系确实存在。

\textbf{E-value 与 Fisher $p$ 值的对比}：在随机化实验中，Fisher $p$ 值衡量的是因果证据。但在大样本观测性研究中，$p$ 值可能是因果证据的误导性度量——即使真实因果效应为零，微弱的未测混杂也可以产生极小的 $p$ 值。采样误差随样本量增大而减小，但未测混杂带来的不确定性\textbf{不会}随样本量增大而消失。E-value 正是为了填补这一空白而设计的。

\begin{Lemma}[$\beta$ 函数性质]\label{lm:17-beta-properties}
定义 $\beta(w_1, w_2) = w_1 w_2 / (w_1 + w_2 - 1)$（$w_1 > 1, w_2 > 1$）。则：
\begin{enumerate}
  \item $\beta(w_1, w_2)$ 对称于 $w_1$ 和 $w_2$；
  \item $\beta(w_1, w_2)$ 对 $w_1$ 和 $w_2$ 均单调递增；
  \item $\beta(w_1, w_2) \leq w_1$ 且 $\beta(w_1, w_2) \leq w_2$；
  \item $\beta(w_1, w_2) \leq w^2 / (2w - 1)$，其中 $w = \max(w_1, w_2)$。
\end{enumerate}
\end{Lemma}

\textbf{引理 \ref{lm:17-beta-properties} 的意义}：(3) 意味着 Cornfield 不等式可以直接从定理 \ref{th:17-cornfield} 推出；(4) 是连接 Cornfield 界与 E-value 的关键步骤。\footnote{证明见附录 \cref{sec:appendix-17-lemma-proof}}

\section{经典例子：吸烟与肺癌}\label{sec:17-smoking-lung-cancer}

\begin{Example}[Hammond and Horn, 1958]\label{ex:17-smoking-lung-cancer}
研究美国人群中吸烟与肺癌的关系。不调整协变量的数据可列为下表：

\begin{center}
\begin{tabular}{lcc}
\toprule
           & 肺癌 & 无肺癌 \\ \midrule
吸烟者      & 397  & 78,557 \\
非吸烟者    & 51   & 108,778 \\ \bottomrule
\end{tabular}
\end{center}

基于该数据，风险比估计为 $\widehat{\RR} = 10.73$，95\% 置信区间为 $[8.02, 14.36]$。

\begin{itemize}
  \item 解释点估计的 E-value：
    \[
    \mathrm{E\text{-}value} = 10.73 + \sqrt{10.73 \times 9.73} \approx 20.95 .
    \]
  \item 解释下限的 E-value：
    \[
    \mathrm{E\text{-}value}_{\mathrm{CI}} = 8.02 + \sqrt{8.02 \times 7.02} \approx 15.52 .
    \]
\end{itemize}

\textbf{解读}：要完全解释掉观察到的风险比（甚至其置信下限），未测混杂与治疗的关联以及与结果的关联，都必须达到至少 15.52。这意味着未测混杂必须是一个\textbf{极强}的风险因子——它与吸烟的关联程度要接近吸烟本身与肺癌的关联程度（10 倍以上）。在领域知识上，这种强度的未知混杂似乎不太合理，因此该研究提供了较强的因果证据。
\end{Example}

图 \ref{fig:17-confounding-contour}（Hammond and Horn 数据的混淆强度轮廓）展示了要解释掉点估计和置信下限，两类混淆强度 $(\RR_{ZU\mid x}, \RR_{UY\mid x})$ 需要满足的条件。

\begin{figure}[ht]
\centering
% Placeholder for confounding intensity contour figure
\fbox{\begin{minipage}{0.6\textwidth}
\vspace{2cm}
\centering
混淆强度轮廓图\\
(待插入)
\vspace{2cm}
\end{minipage}}
\caption{Hammond and Horn 数据的混淆强度轮廓：$\RR_{ZU\mid x}$（$x$ 轴）与 $\RR_{UY\mid x}$（$y$ 轴）满足 $\RR_{ZY\mid x}^{\mathrm{obs}} \leq \frac{\RR_{ZU\mid x} \cdot \RR_{UY\mid x}}{\RR_{ZU\mid x} + \RR_{UY\mid x} - 1}$ 的区域。}
\label{fig:17-confounding-contour}
\end{figure}

\section{扩展}\label{sec:17-extensions}

\subsection{E-value 与 Bradford Hill 因果判定准则}\label{sec:17-bradford-hill}

E-value 提供了因果证据的定量度量，但它与流行病学中经典的\textbf{Bradford Hill 准则}（Bradford Hill, 1965）有深刻的联系。Bradford Hill 提出了九条判断因果关系的准则，其中第一条是：

\begin{enumerate}
  \item[\textbf{强度（Strength）}]：关联越强，因果解释的可能性越大。
\end{enumerate}

定理 \ref{th:17-cornfield} 从数学上精确化了这一直觉：\textbf{更强的关联需要更强的未测混杂来解释}。E-value 正是这种"强度"的量化。

当然，关联强度只是因果判定的一个方面。Bradford Hill 的其他准则——时间顺序（temporality）、生物学合理性（biological plausibility）、一致性（consistency）、剂量-反应关系（dose-response）——都从不同角度提供证据。E-value 只量化了"强度"这一维度。

\begin{Definition}[Bradford Hill 九条因果准则]\label{def:17-bradford-hill}
\citep{BradfordHill1965} 提出的九条因果判定准则：
\begin{enumerate}
  \item 强度（strength）；
  \item 一致性（consistency）；
  \item 特异性（specificity）；
  \item 时间顺序（temporality）；
  \item 生物学梯度（biological gradient）；
  \item 生物学合理性（plausibility）；
  \item 一致性（coherence）；
  \item 实验证据（experiment）；
  \item 类比（analogy）。
\end{enumerate}
\end{Definition}

E-value 量化了第 (1) 条准则，但没有涉及其他八条。这是 E-value 的适用范围——它是一个有意义的补充，但不是因果判定的全部。

\subsection{逻辑斯蒂回归后的 E-value}\label{sec:17-evalue-logistic}

流行病学研究中，二值结果常用逻辑斯蒂回归建模：
\begin{equation}
\pr(Y_i = 1 \mid Z_i, X_i) = \frac{e^{\beta_0 + \beta_1 Z_i + \beta_2^\top X_i}}{1 + e^{\beta_0 + \beta_1 Z_i + \beta_2^\top X_i}} .
\label{eq:17-logistic-model}
\end{equation}

其中系数 $\beta_1$ 是给定 $X$ 下治疗与结果的\textbf{条件比数比（conditional odds ratio）}的对数。逻辑斯蒂模型的一个重要假设是：比数比在所有 $X$ 取值上是共同的（common odds ratio）。

当结果罕见时（$\pr(Y=1 \mid Z, X)$ 接近 0），条件比数比近似于条件风险比（见附录 \cref{sec:appendix-17-rare-approx}）。

因此，在罕见结果假设下，我们可以直接基于逻辑斯蒂回归估计的 $\hat\beta_1$ 计算 E-value：
\begin{equation}
\widehat{\mathrm{E\text{-}value}} = \exp(\hat\beta_1) + \sqrt{\exp(\hat\beta_1)(\exp(\hat\beta_1) - 1)} .
\label{eq:17-evalue-logistic}
\end{equation}

这使得 E-value 可以方便地应用于常规的流行病学分析流程。

\begin{Example}[NCHS 2003 出生数据]\label{ex:17-nchs}
使用美国国家卫生统计中心 2003 年出生证明数据，研究高龄母亲（$\geq 35$ 岁）对早产的影响。调整社会人口学协变量后，逻辑斯蒂回归得到的条件风险比为 $1.31$，其 95\% 置信区间为 $[1.29, 1.33]$。

对应的 E-value 约为 $1.94$（点估计）和 $1.91$（置信下限）。这意味着要解释掉观察到的关联，未测混杂必须同时与治疗（高龄）和结果（早产）有至少 1.9 倍的关联——这个强度在流行病学上看起来是较大的，但远不如吸烟-肺癌例子中的 20.95。
\end{Example}

\subsection{允许非零真实因果效应}\label{sec:17-nonzero-effect}

定理 \ref{th:17-cornfield} 假设真实因果效应为零。Ding and VanderWeele (2016) 将结果推广到允许非零真实因果效应的情形。

\begin{Theorem}[含真实因果效应的下界]\label{th:17-true-rr-bound}
修改 $\RR_{UY\mid x}$ 的定义为：
\begin{equation}
\RR_{UY\mid x} = \max_{z=0,1} \frac{\pr(Y=1 \mid Z=z, U=1, X=x)}{\pr(Y=1 \mid Z=z, U=0, X=x)} .
\label{eq:17-uy-max}
\end{equation}
则在假设 \cref{eq:17-latent-ignorability} 下，
\begin{equation}
\RR_{ZY\mid x}^{\mathrm{true}} \geq \frac{\RR_{ZY\mid x}^{\mathrm{obs}}}{\frac{\RR_{ZU\mid x} \cdot \RR_{UY\mid x}}{\RR_{ZU\mid x} + \RR_{UY\mid x} - 1}} .
\label{eq:17-true-rr-lower-bound}
\end{equation}
\end{Theorem}

\textbf{定理 \ref{th:17-true-rr-bound} 的解读}：即使存在未测混杂，定理也给出了真实因果效应 $\RR_{ZY\mid x}^{\mathrm{true}}$ 的一个下界。这个下界等于观测风险比除以 Cornfield 上界。当灵敏度参数取不同值时，我们可以评估因果效应下界对未测混杂强度的敏感性。

\section{批评与回应}\label{sec:17-critiques}

自 E-value 论文发表以来，它已成为流行病学研究中报告的标准数字。但它也受到了批评（Ioannidis et al., 2019）。下面讨论几个主要局限性。

\subsection{E-value 不过是风险比的单调变换}\label{sec:17-critique-monotone}

当风险比较大时，E-value 近似等于 $2 \cdot \RR_{ZY\mid x}^{\mathrm{obs}}$——几乎是线性的。因此批评者认为 E-value 不过是风险比的一个单调变换，没有提供额外信息。

这\textbf{部分正确}——E-value 确实完全基于风险比的点估计和置信限。但关键在于：E-value 有\textbf{有意义的解释}（定义 \ref{def:17-evalue}），而风险比本身没有这种"需要多强的未测混杂"的解释。从风险比 10.73 到 E-value 20.95，数字变了，但解释的含义完全不同。

\subsection{E-value 的校准问题}\label{sec:17-critique-calibration}

E-value 等于 $\max(\RR_{ZU\mid x}, \RR_{UY\mid x})$——即要解释掉观察到的关联，未测混杂与治疗和结果关联所必须达到的最小值。问题在于：这个未测混杂本身是无法观测的，因此很难判断某个 E-value 究竟算"大"还是"小"。

此外，E-value 的大小取决于我们调整了多少协变量 $X$——它量化的是\textbf{给定已调整协变量后}的残余混杂强度。不同研究之间的 E-value 不能直接比较，因为它们的 $X$ 不同。

一个直观的校准思路是"逐一剔除协变量"法：如果我们把某个已调整的协变量 $X_j$ 当作"未测混杂"来对待，计算 $Z\text{-}X_j$ 和 $X_j\text{-}Y$ 的风险比，这些比值可以提供参考尺度。但这一方法尚无严格的理论支撑。

\subsection{E-value 最适用于二值结果和风险比}\label{sec:17-critique-binary-only}

定理 \ref{th:17-cornfield} 在二值结果和风险比尺度上表现最佳。Ding and VanderWeele (2016) 也提出了风险差、均值比（针对非负结果）以及生存分析中的风险比（HR）的灵敏度分析方法，但它们不如二值结果下的 E-value 那样优雅和简洁。第 \ref{ch:18} 章将给出平均因果效应的更一般灵敏度分析方法，适用于 outcome regression、IPW 和双重稳健估计。

\section{本章小结}\label{sec:17-summary}

本章围绕 E-value 展开，核心思想是：

\begin{enumerate}
  \item \textbf{未测混杂无法从数据中验证}，但我们可以量化：要解释掉观察到的关联，未测混杂必须有多强？
  \item \textbf{Cornfield 不等式}：要完全解释掉风险比 $\RR_{ZY\mid x}^{\mathrm{obs}}$，未测混杂与治疗的关联和与结果的关联都必须至少达到 $\RR_{ZY\mid x}^{\mathrm{obs}}$；
  \item \textbf{E-value} 是这一思想的精确数学化：$\mathrm{E\text{-}value} = \RR_{ZY\mid x}^{\mathrm{obs}} + \sqrt{\RR_{ZY\mid x}^{\mathrm{obs}}(\RR_{ZY\mid x}^{\mathrm{obs}} - 1)}$；
  \item E-value 越大，意味着需要越强的未测混杂才能解释掉观察到的关联，因果证据越强；
  \item E-value 有明确的局限性——它只量化了 Bradford Hill 准则中的"强度"维度，且最适用于二值结果和风险比。
\end{enumerate}

本章将所有推导放在附录，正文专注于概念的动机、历史的脉络和结果的理解。

\section{习题}\label{sec:17-homework}

\begin{HomeworkProblem}[17.1]\label{hw:17-1}
证明函数 $f(x) = (ax+1)/(bx+1)$ 在 $a > b$ 时关于 $x$ 单调递增，在 $a < b$ 时单调递减。
\end{HomeworkProblem}

\begin{HomeworkProblem}[17.2]\label{hw:17-2}
证明在假设 $Z \perp Y \mid (X,U)$ 下，如果 $\RR_{ZU\mid x} > 1$ 但 $\RR_{UY\mid x} < 1$，则 $\RR_{ZY\mid x}^{\mathrm{obs}} < 1$。这说明 \cref{eq:17-assumptions} 中的第三个条件在某种意义上是冗余的。
\end{HomeworkProblem}

\begin{HomeworkProblem}[17.3]\label{hw:17-3}
证明引理 \ref{lm:17-beta-properties}。
\end{HomeworkProblem}

\begin{HomeworkProblem}[17.4]\label{hw:17-4}
（Schlesselman 1978 公式）在二值治疗 $Z$、结果 $Y$ 和未测混杂 $U$ 下，假设层内（$U=0$ 和 $U=1$）治疗对结果的风险比相同，$\RR_{ZY\mid U=0} = \RR_{ZY\mid U=1}$，且混杂对结果的风险比在两组治疗下也相同（记为 $\gamma$）。证明：
\begin{equation}
\frac{\RR_{ZY}^{\mathrm{obs}}}{\RR_{ZY}^{\mathrm{true}}} = \frac{1 + (\gamma - 1)\pr(U=1 \mid Z=1)}{1 + (\gamma - 1)\pr(U=1 \mid Z=0)} .
\label{eq:17-schlesselman}
\end{equation}
并讨论此公式与定理 \ref{th:17-cornfield} 的关系。
\end{HomeworkProblem}

\begin{HomeworkProblem}[17.5]\label{hw:17-5}
使用与例 \ref{ex:17-nchs} 相同的数据集，计算结果为子痫前期（preeclampsia）时的 E-value。
\end{HomeworkProblem}

\begin{HomeworkProblem}[17.6]\label{hw:17-6}
（风险差的 Cornfield 不等式）考虑二值 $Z, Y, U$。在潜在可忽略性 $Z \perp Y \mid U$ 下，证明风险差的 Cornfield 不等式：
\begin{equation}
\RD_{ZY}^{\mathrm{obs}} = \RD_{ZU} \times \RD_{UY} ,
\label{eq:17-rd-cornfield}
\end{equation}
并由此推出：
\begin{equation}
\min(\RD_{ZU}, \RD_{UY}) \geq \RD_{ZY}^{\mathrm{obs}}, \quad
\max(\RD_{ZU}, \RD_{UY}) \geq \sqrt{\RD_{ZY}^{\mathrm{obs}}} .
\label{eq:17-rd-cornfield-ineq}
\end{equation}
\end{HomeworkProblem}

\section{推荐阅读}\label{sec:17-reading}

\begin{itemize}
  \item \citet{DingVanderWeele2016}：扩展并统一了 Cornfield 型灵敏度分析，是 E-value 的理论基础。
  \item \citet{VanderWeeleDing2017}：首次正式提出 E-value 并给出其流行病学解释。
  \item \citet{BradfordHill1965}：Bradford Hill 原文，九条因果准则的经典来源。
  \item \citet{Ioannidis2019}：对 E-value 的批评性讨论。
\end{itemize}

% % % % % % % % % % % % % % % % % % % % % % % % % % % % % % % % % % % %
% APPENDIX
% % % % % % % % % % % % % % % % % % % % % % % % % % % % % % % % % % % %

